\section{阈值与医疗需求}
\label{sec:sup:medical-link}
参照 Angulo 等\cite{Angulo2021}，设任一时刻需要重症监护的人数约为两个感染仓室人数的固定比例，即 $p_{\rm ICU}(I+I_q)$，并要求它不超过可用的重症监护床位数 $C_{\rm ICU}$。阈值策略只约束 $I$，下面的引理把 $I+I_q$ 与 $\eta$ 联系起来。记
\[
\kappa_I=1+\frac{\beta c_0S_0/N-\gamma}{\delta_q}.
\]

\begin{lemma}
\label{lem:sup:infected-bound}
设条件\eqref{eq:model:trigger}成立，$I_q(0)\le(\kappa_I-1)\eta$。则对 $\mathcal A(\eta)$ 中任一控制和所有 $t\ge0$，有
\[
I(t)+I_q(t)\le\kappa_I\eta .
\]
\end{lemma}
\begin{proof}
$\mathcal A(\eta)$ 中始终有 $I\le\eta$（退出后由定理~\ref{thm:s1:tend}）、$c\le c_0$ 与 $S\le S_0$，故新增感染 $F=\beta cSI/N\le\beta c_0S_0\eta/N$。阈值控制期内，由式\eqref{eq:model:balance}，进入 $I_q$ 的流量为 $F-\gamma\eta$；其余时段控制为 $(c_0,q_0)$，流量为 $q_0F$，而 $S_0>S_c$ 给出 $q_0\beta c_0S_0/N\le\beta c_0S_0/N-\gamma$。两种情形下流量都不超过 $\delta_q(\kappa_I-1)\eta$，于是几乎处处有
\[
\dot I_q\le\delta_q\bigl[(\kappa_I-1)\eta-I_q\bigr],
\]
由 $I_q(0)\le(\kappa_I-1)\eta$ 即得 $I_q\le(\kappa_I-1)\eta$。再与 $I\le\eta$ 相加。
\end{proof}

因此取
\[
\eta\le\frac{C_{\rm ICU}}{p_{\rm ICU}\kappa_I}
\]
时，阈值策略的各控制都满足 $p_{\rm ICU}(I+I_q)\le C_{\rm ICU}$。$\kappa_I-1$ 是每名非隔离感染者对应的隔离感染者人数的上界：一名非隔离感染者在常规接触下最多引起 $\beta c_0S_0/(\gamma N)$ 例新感染，其中一例补充 $I$，其余被追踪进入 $I_q$。正文表~\ref{tab:joint:compare}与表~\ref{tab:xian:joint-compare}末列的 $\max_t(I+I_q)/\eta$ 即名义重症监护需求峰值与 $p_{\rm ICU}\eta$ 之比。

这一比例近似与 Angulo 等的处理相同，没有考虑入院延迟与重症监护时长。基准算例的阈值 $0.05N$ 只用于演示机制，不由这一条件换算；西安情景的取值见第~\ref{sec:sup:threshold-scale}~节。
